% systems-biology-shaders.tex
% Systems Biology Shaders: A GPU-Native Formalism for Cellular
% Observation Through the Rendering-Measurement Identity
%
% Compile with:
%   pdflatex systems-biology-shaders
%   bibtex systems-biology-shaders
%   pdflatex systems-biology-shaders
%   pdflatex systems-biology-shaders

\documentclass[11pt,twocolumn,a4paper]{article}

% ─── Packages ────────────────────────────────────────────────────────
\usepackage[utf8]{inputenc}
\usepackage[T1]{fontenc}
\usepackage{lmodern}
\usepackage{amsmath,amssymb,amsthm,mathtools}
\usepackage{physics}
\usepackage{bm}
\usepackage{graphicx}
\usepackage{xcolor}
\usepackage{booktabs}
\usepackage{array}
\usepackage{multirow}
\usepackage{enumitem}
\usepackage{float}
\usepackage{caption}
\usepackage{subcaption}
\usepackage[colorlinks=true,linkcolor=blue!60!black,citecolor=blue!60!black,
            urlcolor=blue!60!black]{hyperref}
\usepackage[margin=1.8cm,top=2.2cm,bottom=2.2cm]{geometry}
\usepackage{natbib}
\usepackage{fancyhdr}
\usepackage{titlesec}
\usepackage{microtype}
\usepackage{algorithm}
\usepackage{algpseudocode}
\usepackage{listings}
\usepackage{tabularx}

% ─── Listings Style ─────────────────────────────────────────────────
\lstdefinestyle{glsl}{
  language=C,
  basicstyle=\ttfamily\scriptsize,
  keywordstyle=\color{blue!70!black}\bfseries,
  commentstyle=\color{green!50!black}\itshape,
  stringstyle=\color{red!60!black},
  numbers=left,
  numberstyle=\tiny\color{gray},
  numbersep=4pt,
  frame=single,
  breaklines=true,
  columns=fullflexible,
  morekeywords={vec2,vec3,vec4,mat2,mat3,mat4,sampler2D,sampler3D,
    uniform,varying,in,out,layout,location,precision,
    highp,mediump,lowp,texture,gl_FragColor,void,float,int,
    gl_VertexID,gl_Position}
}

% ─── Theorem Environments ────────────────────────────────────────────
\newtheorem{theorem}{Theorem}[section]
\newtheorem{lemma}[theorem]{Lemma}
\newtheorem{corollary}[theorem]{Corollary}
\newtheorem{proposition}[theorem]{Proposition}
\newtheorem{definition}[theorem]{Definition}
\newtheorem{remark}[theorem]{Remark}
\newtheorem{axiom}{Axiom}
\newtheorem{example}[theorem]{Example}

% ─── Custom Commands ─────────────────────────────────────────────────
\newcommand{\kb}{k_{\mathrm{B}}}
\newcommand{\Sk}{S_k}
\newcommand{\St}{S_t}
\newcommand{\Se}{S_e}
\newcommand{\BigO}{\mathcal{O}}
\newcommand{\R}{\mathbb{R}}
\newcommand{\N}{\mathbb{N}}
\newcommand{\Z}{\mathbb{Z}}
\newcommand{\Hcat}{\mathcal{H}}
\newcommand{\dcat}{d_{\mathrm{cat}}}
\newcommand{\Gcond}{\mathcal{G}}
\newcommand{\Back}{\mathcal{B}}
\newcommand{\Obs}{\mathcal{O}}
\newcommand{\Rend}{\mathcal{R}}
\newcommand{\Meas}{\mathcal{M}}
\newcommand{\Tex}{\mathcal{T}}
\newcommand{\Part}{\mathcal{P}}
\newcommand{\Cat}{\mathcal{C}}
\newcommand{\taup}{\tau_{\mathrm{p}}}
\newcommand{\Mgpu}{M_{\mathrm{GPU}}}
\newcommand{\Mshader}{M_{\mathrm{shader}}}
\newcommand{\DP}{\Delta P}
\newcommand{\Tref}{T_{\mathrm{ref}}}
\newcommand{\trec}{t_{\mathrm{rec}}}
\newcommand{\fref}{f_{\mathrm{ref}}}
\DeclareMathOperator{\Tr}{Tr}
\DeclareMathOperator{\diag}{diag}
\DeclareMathOperator{\sgn}{sgn}

% ─── Page Style ──────────────────────────────────────────────────────
\pagestyle{fancy}
\fancyhf{}
\fancyhead[L]{\small\textit{Systems Biology Shaders}}
\fancyhead[R]{\small\thepage}
\renewcommand{\headrulewidth}{0.4pt}

% ─── Title Section Formatting ────────────────────────────────────────
\titleformat{\section}{\large\bfseries}{\thesection.}{0.5em}{}
\titleformat{\subsection}{\normalsize\bfseries}{\thesubsection.}{0.5em}{}
\titleformat{\subsubsection}{\normalsize\itshape}{\thesubsubsection.}{0.5em}{}

% ═══════════════════════════════════════════════════════════════════════
\begin{document}

\twocolumn[
\begin{@twocolumnfalse}
\begin{center}

{\LARGE\bfseries Systems Biology Shaders:\\[4pt]
A GPU-Native Formalism for Cellular Observation\\[4pt]
Through the Rendering--Measurement Identity}

\vspace{12pt}

{\large Kundai Farai Sachikonye}

\vspace{6pt}

{\small\itshape Technical University of Munich, School of Life Sciences\\[0.3em]
Chair of Proteomics and Bioanalytics\\[0.3em]
\texttt{kundai.sachikonye@wzw.tum.de}}

\vspace{6pt}

{\small May 2026}

\vspace{14pt}

\begin{minipage}{0.92\textwidth}
\textbf{Abstract.}
Systems biology has standardised the \emph{description} of biological
models through markup languages (SBML, CellML, BioPAX), yet the
computational paradigm remains forward numerical simulation---a process
we prove to be informationally equivalent to the cell itself and
therefore incapable of producing knowledge the cell does not already
possess.  We introduce \emph{Systems Biology Shaders} (SBS), a
formalism in which biological systems are specified not as markup but
as GPU fragment shaders whose rendering constitutes physical
measurement.  The theoretical foundation rests on three results derived
from first principles.  First, we prove the \emph{Oscillatory Necessity
Theorem}: every persistent dynamical system confined to bounded phase
space necessarily oscillates, eliminating static, monotonic, and
non-recurrent modes.  Second, we prove the \emph{Triple Equivalence
Theorem}: the oscillatory, categorical, and partitional descriptions
of any bounded system yield identical state counts, identical
entropies, and are connected by explicit bijections---establishing
mathematical identity, not analogy.  Third, we prove the
\emph{Rendering--Measurement Identity}: for a fragment shader
implementing partition observation, the act of rendering a texture is
identical to the act of measuring the categorical state; the texture
is the measurement result itself, not a depiction of it.  From these
results we derive: (i)~a cellular circuit model in which molecular
species are nodes carrying chemical potential, reactions are edges
with conductance, and Kirchhoff's laws hold exactly; (ii)~a
time-invariant backward trajectory that serves as the categorical
address of each molecular node, enabling observation to proceed via
dimensional channels rather than temporal evolution; (iii)~an
$\BigO(1)$-memory streaming protocol requiring $\sim\!13$\,MB
independent of database size; (iv)~a \emph{Composition Inflation
Theorem} showing that $n$ observation events across $d$ channels yield
$T(n,d) = d\,(1+d)^{n-1}$ distinguishable trajectories; and (v)~a
complete GLSL shader pipeline (vertex, observation, interference,
quality metric, and display shaders) that implements the formalism in
any WebGL\,2 browser.  We specify the SBS format, provide reference
shader implementations, and demonstrate applications to cell-state
instantiation from partial observations, disease detection via loop
holonomy, and therapeutic effect prediction.  SBS replaces the
markup-then-simulate paradigm with a shader-is-the-model paradigm:
the biological model is the shader program, the observation is the
render call, and the result is the texture.

\vspace{8pt}
\noindent\textbf{Keywords:} systems biology shaders, fragment shader,
rendering--measurement identity, cellular circuit model, oscillatory
necessity, triple equivalence, partition coordinates, S-entropy,
time-invariant trajectory, GPU observation, SBML alternative,
composition inflation, loop holonomy
\end{minipage}

\vspace{18pt}
\end{center}
\end{@twocolumnfalse}
]

% ═══════════════════════════════════════════════════════════════════════
%                    PART I: THE OBSERVATION PROBLEM
% ═══════════════════════════════════════════════════════════════════════

\section{Introduction}
\label{sec:introduction}

\subsection{The Markup Paradigm in Systems Biology}

The dominant representational paradigm in computational systems biology
is \emph{markup}: a biological model is encoded as a structured
document describing species, reactions, compartments, and kinetic
laws.  The Systems Biology Markup Language
(SBML)~\citep{hucka2003sbml,keating2020sbml} is the most widely
adopted standard, with CellML~\citep{lloyd2004cellml} and
BioPAX~\citep{demir2010biopax} serving complementary roles for
electrophysiology and pathway exchange respectively.

These standards have been enormously successful at enabling model
sharing, reproducibility, and tool interoperability.  Yet they share a
common limitation: they describe \emph{what} a biological system is
(its components and their relationships) but delegate \emph{what to do
with it} to a separate numerical integrator.  The computational
pipeline is:
\begin{equation}
  \text{Markup} \;\xrightarrow{\text{parse}}\;
  \text{ODE system} \;\xrightarrow{\text{integrate}}\;
  \text{Trajectory } \mathbf{x}(t).
  \label{eq:markup-pipeline}
\end{equation}
The integrator---typically CVODE, LSODA, or the Gillespie
algorithm~\citep{gillespie1977}---steps the system forward in time from
an initial condition, producing a trajectory that is then
\emph{visualised} on a display.

This pipeline contains a deep epistemological flaw that has gone
largely unexamined.

\subsection{The Simulation Problem}

Consider what forward simulation actually computes.  Given the state
$\mathbf{x}_0 \in \R^N$ of $N$ molecular species and the
stoichiometric system
\begin{equation}
  \dot{\mathbf{x}} = \mathbf{S}\,\mathbf{v}(\mathbf{x}),
  \label{eq:ode}
\end{equation}
where $\mathbf{S}$ is the $N \times R$ stoichiometric matrix and
$\mathbf{v} : \R^N \to \R^R$ is the rate vector, the integrator
produces $\mathbf{x}(t)$ for $t > 0$.  But this trajectory is, by
construction, a faithful representation of the system's own dynamics.
It does not transcend those dynamics; it embodies them.

\begin{proposition}[Informational Equivalence of Simulation]
\label{prop:sim-equiv}
Let $\Omega$ denote the set of trajectories consistent with
\eqref{eq:ode}.  A forward simulation starting from $\mathbf{x}_0$
selects a single $\omega \in \Omega$ (or a distribution over $\Omega$
in the stochastic case).  The information content of this selection
equals that of the actual cellular trajectory, because both are
governed by the same probability measure on $\Omega$.  The simulation
produces zero surplus information.
\end{proposition}

\begin{proof}
Let $P$ be the probability measure on $\Omega$ induced by the kinetics
and initial condition.  The simulator samples $\omega \sim P$.  The
cell generates $\omega' \sim P$ by the same kinetics.  Both provide
i.i.d.\ samples from $P$.  The mutual information between the
simulation output and any property not already determined by $P$ is
zero:
$I(\omega; Q \mid P) = 0$ for any query $Q$ not a measurable function
of $P$.  \qed
\end{proof}

The implication is immediate: if forward simulation could reveal
disease, cells would self-diagnose.  They cannot.

\begin{definition}[Epistemic Blindness]
\label{def:blind}
A dynamical system $(X, \mathbf{v})$ is \emph{epistemically blind}
with respect to a partition $\{A, B\}$ of $X$ if the flow
$\mathbf{v}(x)$ does not encode membership in $A$ or $B$: there
exists no function $f : TX \to \{A, B\}$ such that $f(\mathbf{v}(x))$
correctly classifies $x$ for all $x \in X$.
\end{definition}

\begin{proposition}[Cells Are Epistemically Blind to Disease]
\label{prop:blind}
Under generic kinetic models (mass-action, Michaelis--Menten), a cell
operating via forward reaction dynamics is epistemically blind with
respect to the partition \emph{healthy}/\emph{diseased}.
\end{proposition}

\begin{proof}
The rate law $v_{ij}([C_i], [C_j], k_{ij})$ depends only on current
concentrations and rate constants.  It carries no reference to a
healthy state $\mathbf{x}_H^*$.  Disease states frequently involve
regulatory rewiring or epigenetic changes that alter the parameters of
$\mathbf{v}$ rather than concentrations directly, making the current
flux indistinguishable from that of a normal cell elsewhere in the
healthy attractor.  No function of $\mathbf{v}(\mathbf{x}_D^*)$ alone
reliably distinguishes diseased from healthy.  \qed
\end{proof}

\subsection{From Markup to Shaders}

The resolution lies in abandoning forward simulation in favour of
\emph{observation}: determining the complete state from partial
measurements by exploiting the topological and thermodynamic structure
of the reaction network.  This requires an external observer who knows
the network topology and the constraint structure.

We propose that this observer is a GPU fragment shader.

The central claim of this paper is that a fragment shader is not a
visualisation tool.  When it evaluates a function at a texel coordinate
and writes the result to a framebuffer, it is not creating a picture
of a measurement---it \emph{is} performing the measurement.  The
rendered texture is the computed result in categorical representation,
not a depiction of it.  This is not metaphor.  We prove it as a
mathematical identity (Theorem~\ref{thm:rmi}).

We call this formalism \emph{Systems Biology Shaders} (SBS).  Where
SBML specifies a biological model as markup to be parsed and
simulated, SBS specifies a biological model as a shader program to be
rendered.  The model is the shader.  The observation is the render
call.  The result is the texture.

\subsection{Contributions}

\begin{enumerate}[leftmargin=*,itemsep=2pt]
\item We prove the \emph{Oscillatory Necessity Theorem}: every
  persistent bounded system must oscillate
  (\S\ref{sec:oscillatory-necessity}).

\item We prove the \emph{Triple Equivalence Theorem}: oscillation,
  categorical distinction, and partition operation are mathematically
  identical (\S\ref{sec:triple-equivalence}).

\item We derive a \emph{cellular circuit model} from first principles
  of thermodynamics, with exact Kirchhoff analogs
  (\S\ref{sec:circuit-model}).

\item We prove the \emph{Rendering--Measurement Identity}: fragment
  shader evaluation is physical measurement
  (\S\ref{sec:rendering-measurement}).

\item We prove that multi-channel observation provides dimensional
  progression without temporal evolution
  (\S\ref{sec:temporal-observation}).

\item We derive the \emph{Composition Inflation Theorem}:
  $T(n,d) = d\,(1+d)^{n-1}$ (\S\ref{sec:composition-inflation}).

\item We specify the complete SBS pipeline with reference GLSL
  implementations (\S\ref{sec:sbs-specification}).

\item We demonstrate applications to disease detection, drug design,
  and cell-state instantiation (\S\ref{sec:applications}).
\end{enumerate}


% ═══════════════════════════════════════════════════════════════════════
%              PART II: THEORETICAL FOUNDATIONS
% ═══════════════════════════════════════════════════════════════════════

\section{Bounded Phase Space and Oscillatory Necessity}
\label{sec:oscillatory-necessity}

\subsection{Axioms}

We begin with axioms that any biological system satisfies by
definition.

\begin{axiom}[Bounded Phase Space]
\label{ax:bounded}
The phase space $X$ of a living cell is a compact subset of $\R^N$:
every concentration, potential, and flux is bounded above and below by
physical constraints (solubility limits, membrane integrity, energy
availability).
\end{axiom}

\begin{axiom}[Persistence]
\label{ax:persistence}
The dynamical system $(X, \mathbf{v})$ persists: the trajectory
$\mathbf{x}(t)$ exists for all $t \geq 0$ and remains within $X$.
\end{axiom}

\begin{axiom}[Reproducibility]
\label{ax:reproducibility}
Given identical initial conditions and parameters, the system produces
statistically identical trajectories.
\end{axiom}

These three axioms---boundedness, persistence, and
reproducibility---are the only assumptions.  Everything that follows
is derived from them.

\subsection{Elimination of Non-Oscillatory Modes}

\begin{lemma}[No Static Mode]
\label{lem:no-static}
A bounded persistent system with $N \geq 2$ interacting species cannot
remain at a static fixed point indefinitely, provided the interaction
graph is connected and at least one edge has nonzero flux.
\end{lemma}

\begin{proof}
At a fixed point $\mathbf{x}^*$, $\mathbf{S}\,\mathbf{v}(\mathbf{x}^*)
= \mathbf{0}$.  For a connected reaction network with at least one
nonzero flux, this requires $\mathbf{v}(\mathbf{x}^*)$ to lie in the
null space of $\mathbf{S}$.  By the rank-nullity theorem,
$\dim(\ker \mathbf{S}) = R - \mathrm{rank}(\mathbf{S})$.  In a
connected network with $N \geq 2$ species and generic (non-degenerate)
rate constants, the stoichiometric compatibility class
$\{\mathbf{x} : \mathbf{x} = \mathbf{x}_0 + \mathbf{S}\boldsymbol{\xi},\;
\boldsymbol{\xi} \geq 0\}$ has dimension
$\mathrm{rank}(\mathbf{S}) \geq 1$.  Perturbations along this class
are amplified by at least one reaction, destabilising the fixed point.
Generic connected networks of $N \geq 2$ species therefore do not
admit globally attracting fixed points.  \qed
\end{proof}

\begin{lemma}[No Monotonic Mode]
\label{lem:no-monotonic}
A bounded persistent system cannot exhibit monotonically increasing or
decreasing trajectories for any coordinate.
\end{lemma}

\begin{proof}
If $x_i(t)$ is monotonically increasing and bounded above, then by
the Monotone Convergence Theorem $x_i(t) \to x_i^* < \infty$.  But
then $\dot{x}_i \to 0$, and $x_i$ ceases to increase, contradicting
monotonicity for all $t$.  The decreasing case is symmetric.  \qed
\end{proof}

\begin{lemma}[No Non-Recurrent Chaos]
\label{lem:no-chaos}
A bounded persistent reproducible system cannot exhibit
non-recurrent chaotic dynamics.
\end{lemma}

\begin{proof}
By the Poincar\'{e} Recurrence Theorem~\citep{poincare1890,katok1995},
any measure-preserving flow on a compact space returns arbitrarily
close to any initial state.  Reproducibility (Axiom~\ref{ax:reproducibility})
requires that the system generates the same statistical distribution
over trajectories, which---combined with compactness---ensures the
existence of a finite invariant measure.  Non-recurrent trajectories
(those that never return to a neighbourhood of their starting point)
have measure zero under this invariant measure and are therefore
eliminated.  \qed
\end{proof}

\begin{theorem}[Oscillatory Necessity]
\label{thm:oscillatory-necessity}
Every persistent dynamical system confined to bounded phase space with
connected interaction graph necessarily exhibits oscillatory dynamics.
\end{theorem}

\begin{proof}
Lemma~\ref{lem:no-static} eliminates static equilibria.
Lemma~\ref{lem:no-monotonic} eliminates monotonic trajectories.
Lemma~\ref{lem:no-chaos} eliminates non-recurrent modes.  The
remaining dynamics must be recurrent and non-monotonic on a compact
set, which---by the Jordan Curve Theorem in the planar case and the
Poincar\'{e}--Bendixson classification in general---implies
oscillatory (periodic or quasi-periodic) motion.  \qed
\end{proof}

\begin{corollary}[Biological Systems Oscillate]
Living cells satisfy Axioms~\ref{ax:bounded}--\ref{ax:reproducibility}
and have connected metabolic networks.  Therefore, all persistent
cellular dynamics are oscillatory.  This is not a modelling assumption
but a mathematical consequence of boundedness and persistence.
\end{corollary}

This result is consistent with the extensive experimental literature
on biological oscillations~\citep{goldbeter1996,novak2008,winfree1967,
tyson2003}.


\section{Partition Coordinates and Shell Capacity}
\label{sec:partition-coordinates}

\subsection{Hierarchical Partition of Bounded Phase Space}

The bounded phase space $X \subset \R^N$ admits a hierarchical
partition into shells of increasing depth.

\begin{definition}[Partition Depth]
\label{def:partition-depth}
For a bounded interval $[a, b]$ in one dimension, define the
\emph{partition at depth $n$} as the division of $[a, b]$ into $n$
equal subintervals.  Each subinterval is further characterised by an
angular momentum quantum number $\ell \in \{0, 1, \ldots, n-1\}$
specifying the subpartition, a magnetic quantum number
$m \in \{-\ell, \ldots, +\ell\}$ specifying orientation within the
subpartition, and a binary spin quantum number $s \in \{-\tfrac{1}{2},
+\tfrac{1}{2}\}$ specifying parity.
\end{definition}

\begin{theorem}[Shell Capacity]
\label{thm:shell-capacity}
The number of distinguishable states at partition depth $n$ is
\begin{equation}
  C(n) = 2n^2.
  \label{eq:shell-capacity}
\end{equation}
\end{theorem}

\begin{proof}
At depth $n$, the angular momentum quantum number takes values
$\ell = 0, 1, \ldots, n-1$.  For each $\ell$, the magnetic quantum
number takes $2\ell + 1$ values.  For each $(\ell, m)$ pair, the spin
takes 2 values.  The total is:
\begin{equation}
  C(n) = \sum_{\ell=0}^{n-1} (2\ell + 1) \cdot 2
       = 2 \sum_{\ell=0}^{n-1} (2\ell + 1)
       = 2 \cdot n^2.
  \label{eq:shell-capacity-proof}
\end{equation}
The last equality follows from the identity
$\sum_{\ell=0}^{n-1}(2\ell+1) = n^2$.  \qed
\end{proof}

\begin{remark}
The formula $C(n) = 2n^2$ is identical to the electron shell capacity
of atomic physics~\citep{griffiths2018}.  This is not a coincidence:
both arise from hierarchical partition of a bounded space.  The
partition coordinates $(n, \ell, m, s)$ are not ``quantum numbers
borrowed by analogy''---they are the natural coordinates of any
hierarchically partitioned bounded space.
\end{remark}


\section{The Triple Equivalence Theorem}
\label{sec:triple-equivalence}

We now prove that three descriptions of a bounded system's state
are mathematically identical.

\begin{definition}[Oscillatory Description]
\label{def:oscillatory}
Given $M$ independent bounded degrees of freedom, each oscillating
(by Theorem~\ref{thm:oscillatory-necessity}) with phase
$\phi_i \in [0, 2\pi)$, the oscillatory state is the phase vector
$\boldsymbol{\phi} = (\phi_1, \ldots, \phi_M) \in [0, 2\pi)^M$.
If each phase is resolved to $n$ distinguishable levels, the
oscillatory state count is $\Omega_O = n^M$.
\end{definition}

\begin{definition}[Categorical Description]
\label{def:categorical}
A \emph{categorical distinction}~\citep{maclane1998} on a set $X$ is
a morphism $f : X \to \{1, \ldots, n\}$ assigning each element to one
of $n$ categories.  For $M$ independent objects, the categorical state
is a tuple $(c_1, \ldots, c_M) \in \{1, \ldots, n\}^M$, giving
$\Omega_C = n^M$.
\end{definition}

\begin{definition}[Partition Description]
\label{def:partitional}
A \emph{partition} of a set $X$ into $n$ blocks assigns each element
to exactly one block.  For $M$ independent sets, the partition state
is $(p_1, \ldots, p_M) \in \{1, \ldots, n\}^M$, giving
$\Omega_P = n^M$.
\end{definition}

\begin{theorem}[Triple Equivalence]
\label{thm:triple-equivalence}
For $M$ independent degrees of freedom, each with $n$ distinguishable
states, the oscillatory, categorical, and partition descriptions are
mathematically identical:
\begin{enumerate}[leftmargin=*,itemsep=1pt]
\item \textbf{State counts:}
  $\Omega_O = \Omega_C = \Omega_P = n^M$.
\item \textbf{Entropies:}
  $S_O = S_C = S_P = \kb M \ln n$.
\item \textbf{Bijective maps:} There exist explicit bijections
  $\alpha : [0,2\pi)^M \to \{1,\ldots,n\}^M$,
  $\beta : \{1,\ldots,n\}^M \to \{1,\ldots,n\}^M$, and
  $\gamma : \{1,\ldots,n\}^M \to [0,2\pi)^M$ forming a closed
  commutative triangle $\gamma \circ \beta \circ \alpha = \mathrm{id}$.
\end{enumerate}
\end{theorem}

\begin{proof}
\textbf{State counts and entropies.}  All three descriptions assign
each of $M$ independent degrees of freedom to one of $n$ labels.
The state space cardinality is $n^M$ in each case.  The Boltzmann
entropy $S = \kb \ln \Omega$ gives $S = \kb M \ln n$ for all three.

\textbf{Bijections.}  Define
$\alpha(\phi_1, \ldots, \phi_M) = (\lfloor n\phi_1/(2\pi) \rfloor + 1,
\ldots, \lfloor n\phi_M/(2\pi) \rfloor + 1)$, which maps each
oscillatory phase to its partition bin.  Define $\beta$ as the
identity on $\{1,\ldots,n\}^M$ (categorical and partition labels are
the same set).  Define $\gamma(c_1, \ldots, c_M) = (2\pi(c_1 -
\tfrac{1}{2})/n, \ldots, 2\pi(c_M - \tfrac{1}{2})/n)$, which maps
each categorical label to the midpoint phase of its bin.  Then
$\gamma \circ \beta \circ \alpha$ maps a phase vector to its
bin midpoint, which is the identity up to intra-bin resolution.
At the resolution of $n$ distinguishable states, the maps are exact
bijections on the quotient space $[0,2\pi)^M / {\sim_n}$.  \qed
\end{proof}

\begin{corollary}[Fundamental Identity]
\label{cor:fundamental}
Oscillation $\equiv$ Categorical Distinction $\equiv$
Partition Operation.  These are three notations for the same
mathematical object, not three analogous objects.
\end{corollary}


\section{The Cellular Circuit Model}
\label{sec:circuit-model}

\subsection{Circuit Construction}

\begin{definition}[Cellular Circuit Graph]
\label{def:circuit}
Given a biochemical network with $N$ molecular species
$\{C_1, \ldots, C_N\}$ and $R$ reactions, define a weighted directed
graph $G = (V, E, w)$ where:
\begin{itemize}[leftmargin=*,nosep]
\item Nodes $V = \{C_1, \ldots, C_N\}$, each carrying a \emph{potential}
  equal to its chemical potential:
  \begin{equation}
    \mu_i = \mu_i^0 + RT \ln [C_i],
    \label{eq:chemical-potential}
  \end{equation}
  interpreted information-theoretically as categorical depth
  $\Hcat_i = -\log_2 P(\sigma_i)$, where $P(\sigma_i) = [C_i]/[C_\mathrm{tot}]$.

\item Edges $E$: for each reaction $C_i \to C_j$ with rate constant
  $k_{ij}$, an edge with \emph{conductance}:
  \begin{equation}
    \Gcond_{ij} = \frac{k_{ij}\,[C_i]}{RT}.
    \label{eq:conductance}
  \end{equation}

\item The \emph{current} through edge $(i,j)$ is the reaction flux:
  \begin{equation}
    J_{ij} = \Gcond_{ij}\,(\mu_i - \mu_j).
    \label{eq:current}
  \end{equation}
\end{itemize}
\end{definition}

\subsection{Kirchhoff's Laws from First Principles}

\begin{theorem}[Kirchhoff's Current Law Analog]
\label{thm:kcl}
At steady state, the sum of currents entering any node equals the sum
leaving:
\begin{equation}
  \sum_{j \in \mathcal{N}(i)} J_{ji} = \sum_{j \in \mathcal{N}(i)} J_{ij}
  \quad \forall\, i \in V,
  \label{eq:kcl}
\end{equation}
where $\mathcal{N}(i)$ is the set of neighbours of node $i$.
\end{theorem}

\begin{proof}
This is stoichiometric mass balance.  At steady state,
$\dot{[C_i]} = 0$, so the total production rate equals the total
consumption rate:
$\sum_j k_{ji}[C_j] = \sum_j k_{ij}[C_i]$.  Multiplying both sides
by $1/(RT)$ and using the definition of conductance and current
recovers~\eqref{eq:kcl}.  \qed
\end{proof}

\begin{theorem}[Kirchhoff's Voltage Law Analog]
\label{thm:kvl}
Around any closed loop $\ell = (i_1, i_2, \ldots, i_k, i_1)$ in the
reaction network, the sum of potential differences vanishes:
\begin{equation}
  \sum_{j=1}^{k} (\mu_{i_j} - \mu_{i_{j+1}}) = 0,
  \label{eq:kvl}
\end{equation}
where indices are cyclic: $i_{k+1} = i_1$.
\end{theorem}

\begin{proof}
This is a telescoping sum:
$\sum_{j=1}^{k}(\mu_{i_j} - \mu_{i_{j+1}}) = \mu_{i_1} - \mu_{i_1}
= 0$.  Thermodynamically, this follows from the Wegscheider
condition~\citep{wegscheider1901}: at detailed balance, the product of
forward rate constants around any cycle equals the product of reverse
rate constants, which implies $\sum_\ell \Delta\mu = 0$.  \qed
\end{proof}

\subsection{Fuzzy State Initialisation}

Biological measurements are inherently imprecise.  We model this using
Zadeh fuzzy membership functions~\citep{zadeh1965,zadeh1975}.

\begin{definition}[Fuzzy Node State]
\label{def:fuzzy}
The state of node $C_i$ is a fuzzy number $\tilde{\mu}_i$ with
membership function $\mu_{\tilde{\mu}_i} : \R \to [0,1]$ satisfying:
(i) normality: $\sup_x \mu_{\tilde{\mu}_i}(x) = 1$;
(ii) convexity: $\mu_{\tilde{\mu}_i}(\lambda x + (1-\lambda)y) \geq
\min(\mu_{\tilde{\mu}_i}(x), \mu_{\tilde{\mu}_i}(y))$ for all
$\lambda \in [0,1]$;
(iii) upper semicontinuity.
\end{definition}

The Zadeh extension principle lifts the Kirchhoff
equations~\eqref{eq:kcl}--\eqref{eq:kvl} to fuzzy arithmetic: if
$\tilde{\mu}_i$ and $\tilde{\mu}_j$ are fuzzy, then
$\tilde{J}_{ij} = \tilde{\Gcond}_{ij} \cdot (\tilde{\mu}_i -
\tilde{\mu}_j)$ is computed via $\alpha$-cut
arithmetic~\citep{zadeh1975}.

\subsection{Time-Invariant Backward Trajectories}

\begin{definition}[Backward Trajectory]
\label{def:backward}
The \emph{backward trajectory} of node $C_i$ is the
maximum-a-posteriori (MAP) sequence of predecessor nodes:
\begin{equation}
  \Back(C_i) = \arg\max_{(C_{j_1}, \ldots, C_{j_T})}
  \prod_{t=1}^{T} P(C_{j_{t-1}} \to C_{j_t} \mid G, \mathbf{x}),
  \label{eq:backward}
\end{equation}
computed by the Viterbi algorithm~\citep{viterbi1967,rabiner1989} on
the circuit graph $G$.
\end{definition}

\begin{theorem}[Time Invariance]
\label{thm:time-invariance}
For an autonomous kinetic system (time-independent rate constants),
the MAP backward trajectory $\Back(C_i)$ is time-invariant: it
depends on the current state $\mathbf{x}$ and network topology $G$
but not on the absolute time $t$ at which observation occurs.
\end{theorem}

\begin{proof}
The transition probabilities
$P(C_j \to C_k) = k_{jk}[C_j] / \sum_m k_{jm}[C_j]
= k_{jk} / \sum_m k_{jm}$ depend only on rate constants (which are
time-independent by assumption) and the branching ratios of the
current state.  The MAP sequence is the path maximising the product
of these probabilities, which is a function of $(G, \mathbf{x})$
alone.  \qed
\end{proof}

\begin{corollary}[Categorical Address]
\label{cor:categorical-address}
The backward trajectory $\Back(C_i)$ is a \emph{categorical address}:
a fixed geometric identity of node $C_i$ within the network's state
space.  It serves the same role as a postal address---it locates the
node without reference to time.
\end{corollary}

\subsection{Trajectory Completion}

The circuit model enables \emph{trajectory completion}: determining
the complete state of every node from partial observations by
iterating three constraint classes.

\begin{algorithm}[H]
\caption{Trajectory Completion Algorithm}
\label{alg:completion}
\begin{algorithmic}[1]
\Require Partial observations $\{(\tilde{\mu}_i, i \in S)\}$,
  circuit $G = (V, E, w)$
\Ensure Complete state $\{\hat{\mu}_i, i \in V\}$
\State Initialise unobserved nodes with uniform fuzzy states
\Repeat
  \State \textbf{KCL pass:} Propagate fuzzy currents at each node
  \State \textbf{KVL pass:} Enforce loop consistency
  \State \textbf{Backward pass:} Update MAP trajectories via Viterbi
  \State Defuzzify: update crisp estimates
    $\hat{\mu}_i \gets \mathrm{centroid}(\tilde{\mu}_i)$
\Until{$\max_i |\hat{\mu}_i^{(k)} - \hat{\mu}_i^{(k-1)}| < \epsilon$}
\end{algorithmic}
\end{algorithm}

\begin{theorem}[Convergence]
\label{thm:convergence}
The Trajectory Completion Algorithm converges to a unique fixed point
on the space of fuzzy state tuples equipped with the Hausdorff
metric~\citep{beer1993}, provided the circuit graph $G$ is connected
and the conductances satisfy $\Gcond_{ij} > 0$ for all edges.
\end{theorem}

\begin{proof}
Define the operator $T$ mapping a fuzzy state tuple
$\tilde{\boldsymbol{\mu}}$ to the result of one complete
KCL--KVL--backward iteration.  The fuzzy state space
$\mathcal{F}^N$ (compact, convex fuzzy numbers on $\R$) is a
complete metric space under the Hausdorff
metric~\citep{beer1993,kreyszig1978}.  Each constraint pass is
a contraction: (i)~KCL averaging reduces the spread of fuzzy
membership by a factor bounded by the spectral gap of the graph
Laplacian; (ii)~KVL loop closure eliminates inconsistencies, strictly
reducing the Hausdorff diameter; (iii)~Viterbi backward pass selects
the unique MAP path, a projection onto a finite set.  The composition
of these contractions is itself a contraction with constant
$\rho < 1$.  By the Banach Fixed-Point
Theorem~\citep{granas2003}, $T$ has a unique fixed point
$\tilde{\boldsymbol{\mu}}^*$, and the iteration converges
geometrically: $d_H(T^k\tilde{\boldsymbol{\mu}}_0,
\tilde{\boldsymbol{\mu}}^*) \leq \rho^k\,
d_H(\tilde{\boldsymbol{\mu}}_0,
\tilde{\boldsymbol{\mu}}^*)$.  \qed
\end{proof}


% ═══════════════════════════════════════════════════════════════════════
%          PART III: THE RENDERING--MEASUREMENT IDENTITY
% ═══════════════════════════════════════════════════════════════════════

\section{S-Entropy Coordinates}
\label{sec:s-entropy}

To map arbitrary biological observables onto GPU textures, we require
a universal coordinate system.

\begin{definition}[S-Entropy Triple]
\label{def:s-entropy}
For any bounded observable $q$ with support $[q_{\min}, q_{\max}]$,
define three independent entropy coordinates:
\begin{align}
  \Sk &= -\sum_{i} p_i \log p_i \bigg/ \log n,
    \label{eq:sk} \\
  \St &= 1 - \frac{\tau_{\mathrm{corr}}}{\tau_{\max}},
    \label{eq:st} \\
  \Se &= \frac{d\Sk}{dt} \bigg/ \left(\frac{d\Sk}{dt}\right)_{\max},
    \label{eq:se}
\end{align}
where $\Sk$ is the \emph{knowledge entropy} (normalised Shannon
entropy of the partition distribution), $\St$ is the \emph{temporal
entropy} (normalised decorrelation rate), and $\Se$ is the
\emph{evolution entropy} (normalised rate of entropy change).
\end{definition}

\begin{proposition}[Compactness]
\label{prop:compactness}
The S-entropy coordinates $(\Sk, \St, \Se)$ lie in the unit cube
$[0,1]^3$, which is a compact complete metric space under the
Euclidean metric.
\end{proposition}

\begin{proof}
Each coordinate is a ratio of a bounded quantity to its maximum.  By
construction, $0 \leq \Sk \leq 1$, $0 \leq \St \leq 1$, and
$0 \leq \Se \leq 1$.  The product $[0,1]^3$ is compact by
Tychonoff's theorem and complete under the Euclidean metric.  \qed
\end{proof}

\begin{corollary}[GPU Compatibility]
\label{cor:gpu}
S-entropy coordinates map directly to GPU texture coordinates and
colour channels.  A texel at position $(u, v) \in [0,1]^2$ with value
$(\Sk, \St, \Se) \in [0,1]^3$ stored as RGB requires no
normalisation, no unit conversion, and no floating-point scaling.
\end{corollary}


\section{The Rendering--Measurement Identity}
\label{sec:rendering-measurement}

\subsection{Fragment Shader as Observation Apparatus}

\begin{definition}[Fragment Shader]
\label{def:fragment-shader}
A \emph{fragment shader}~\citep{khronosglsl,khronoswebgl2} is a
function $F : [0,1]^2 \times \Theta \to [0,1]^4$ that maps a texel
coordinate $(u,v) \in [0,1]^2$ and a parameter set
$\Theta$ (uniforms) to an RGBA value.  Evaluation is:
(i)~deterministic (same inputs produce same outputs);
(ii)~parallel (each texel is independent);
(iii)~total (defined for every input).
\end{definition}

\begin{definition}[Partition Observation]
\label{def:observation}
An \emph{observation} of a partition state $\sigma$ is a function
$\Obs : \sigma \to [0,1]^4$ that maps the state to a tuple of real
numbers representing physical observables (amplitude, phase, depth,
confidence).
\end{definition}

\begin{theorem}[Rendering--Measurement Identity]
\label{thm:rmi}
For a fragment shader $F$ that implements the observation operation
$\Obs$ on a partition state $\sigma$ encoded as uniforms:
\begin{equation}
  \Rend(F, \sigma) \equiv \Meas(\Obs, \sigma).
  \label{eq:rmi}
\end{equation}
The texture $\Tex = F(\cdot, \sigma)$ output by the fragment shader is
not a representation of the observation result---it \emph{is} the
observation result.
\end{theorem}

\begin{proof}
Let $\sigma$ be a partition state encoded as shader uniforms
(S-entropy coordinates, domain parameters).  The observation $\Obs$
is defined as a function from $\sigma$ to $[0,1]^4$.  The fragment
shader $F$ implements this same function: for each texel $(u,v)$, it
reads $\sigma$ from uniforms and computes
$F(u, v, \sigma) = \Obs(\sigma)(u,v)$.

The identity rests on three properties:
\begin{enumerate}[leftmargin=*,nosep]
\item \textbf{Determinism:} Both $\Obs$ and $F$ are deterministic
  functions of $\sigma$.  Given the same state, they produce the same
  output.  This is guaranteed by the GPU specification: fragment
  shaders are required to be deterministic for fixed inputs and shader
  code~\citep{khronosglsl}.

\item \textbf{Completeness:} $F$ evaluates at every texel in the
  output texture, covering the entire observation domain $[0,1]^2$.
  No region is unobserved.

\item \textbf{Categorical equivalence:} The output of $F$ is stored
  as a texture $\Tex : [0,1]^2 \to [0,1]^4$, which---by
  Corollary~\ref{cor:gpu}---is already in S-entropy coordinates.
  No post-processing or reinterpretation is required.  The texture
  value at $(u,v)$ \emph{is} the observation value at $(u,v)$.
\end{enumerate}

Since $F$ and $\Obs$ are the same function applied to the same input,
producing the same output stored in the same representation,
$\Rend(F,\sigma) = \Meas(\Obs,\sigma)$ as functions
$[0,1]^2 \to [0,1]^4$.  \qed
\end{proof}

\subsection{Texture--State Isomorphism}

\begin{theorem}[Texture--State Isomorphism]
\label{thm:texture-state}
The observation texture $\Tex$ and the partition state $\sigma$ are
isomorphic as information-bearing objects: $\Tex$ contains exactly the
observable information in $\sigma$, no more and no less.
\end{theorem}

\begin{proof}
\textbf{(No less.)}  $\Tex = \Obs(\sigma)$, and $\Obs$ is defined to
extract all observable information from $\sigma$.  Any observable
quantity not in $\Tex$ is not in $\Obs(\sigma)$ and therefore not
observable by definition.

\textbf{(No more.)}  $F$ reads only the uniforms encoding $\sigma$ and
the texel coordinate $(u,v)$.  It has no access to external data, no
random number generator (determinism), and no memory of prior
invocations (statelessness).  Therefore $\Tex$ contains no information
beyond what is in $\sigma$ and the domain geometry.  \qed
\end{proof}

\begin{corollary}[No Separate Computation Step]
\label{cor:no-separate}
The fragment shader does not ``accelerate'' a computation that could
also be performed on the CPU.  There is no separate computational step
that the shader then visualises.  The shader evaluation \emph{is} the
computation.  To ``observe the cell'' and to ``render the texture''
are the same operation.
\end{corollary}


% ═══════════════════════════════════════════════════════════════════════
%        PART IV: MULTI-CHANNEL OBSERVATION WITHOUT TIME
% ═══════════════════════════════════════════════════════════════════════

\section{Temporal Observation and Dimensional Progression}
\label{sec:temporal-observation}

\subsection{The Time-Invariance Constraint}

Theorem~\ref{thm:time-invariance} established that the backward
trajectory of any circuit node is time-invariant: it depends on state
and topology, not on when observation occurs.  This has a profound
consequence for the shader pipeline.

\begin{corollary}[Single-Shot Observation]
\label{cor:single-shot}
Since the categorical address of each node is time-invariant, a single
render pass (one shader evaluation) captures the complete observable
state.  No time-stepping, no animation loop, and no simulation frames
are required.
\end{corollary}

\subsection{Channels as Dimensional Axes}

If time cannot add information (because the trajectory is
time-invariant), how does the system gain distinguishability?  The
answer is \emph{channels}: independent dimensions of observation on
the same node.

\begin{definition}[Observation Channel]
\label{def:channel}
An \emph{observation channel} $\chi = (s, \Obs_s)$ is a pair of a
source identifier $s$ and an observation function $\Obs_s$ that
extracts a specific physical observable from the partition state.
Multiple channels observe different facets of the same node
simultaneously.
\end{definition}

\begin{definition}[Timing Deviation]
\label{def:timing-deviation}
Let $\mathcal{O} = (\fref, \epsilon, M_0)$ be a reference oscillator
with frequency $\fref$.  For the $k$-th observation pulse with
reception time $\trec(k)$, the \emph{timing deviation} is:
\begin{equation}
  \DP(k) = \Tref(k) - \trec(k) \in \R,
  \label{eq:dp}
\end{equation}
where $\Tref(k) = k / \fref$.
\end{definition}

\begin{definition}[Timing Cell]
\label{def:timing-cell}
A \emph{timing cell} is a Borel-measurable subset
$\Cat \subseteq \R$ of the timing-deviation space with positive
Lebesgue measure.  A \emph{cell partition} $\Part = \{\Cat_1, \ldots,
\Cat_m\}$ covers the timing-deviation space with disjoint,
positive-measure cells.
\end{definition}

The positive-measure requirement is physically essential: point
predicates have Lebesgue measure zero and cannot be robustly detected
in the presence of oscillator jitter~\citep{rudin1987}.

\begin{theorem}[Dimensional Progression Without Temporal Evolution]
\label{thm:dimensional-progression}
Let a circuit node $C_i$ be observed through $d$ independent channels
$\{\chi_1, \ldots, \chi_d\}$, each producing a timing deviation
$\DP_s(k)$.  The observable state space of $C_i$ grows as $m^d$
(where $m$ is the number of timing cells per channel) \emph{without
any temporal evolution of the circuit}.
\end{theorem}

\begin{proof}
By Theorem~\ref{thm:time-invariance}, the categorical address of
$C_i$ is time-invariant: its backward trajectory $\Back(C_i)$ does
not change between observation events.  Each channel $\chi_s$ provides
an independent measurement $\DP_s \in \Cat_{j_s}$, where $j_s$
identifies the timing cell.  The $d$-channel observation assigns $C_i$
a tuple $(j_1, \ldots, j_d) \in \{1, \ldots, m\}^d$.  The
cardinality of this observation space is $m^d$, which grows
exponentially in $d$ without requiring the circuit to evolve in time.
The circuit is observed; it does not step.  \qed
\end{proof}

\begin{corollary}[Triple Observation Identity]
\label{cor:triple-observation}
A fragment shader with four output channels (RGBA) simultaneously
computes four independent observables from a single partition-state
read.  For example:
\begin{align}
  R &= \text{optical absorption (amplitude)}, \notag \\
  G &= \text{phase (oscillatory state)}, \notag \\
  B &= \text{partition depth}, \notag \\
  A &= \text{confidence (fuzzy membership)}.
  \label{eq:rgba-observables}
\end{align}
Four channels on a single texel.  Zero time advancement.
\end{corollary}


\section{The Composition Inflation Theorem}
\label{sec:composition-inflation}

When multiple channels are combined into trajectories, the
distinguishability grows faster than exponentially.

\begin{definition}[Labeled Composition]
\label{def:labeled-composition}
A \emph{$d$-labeled composition of $n$} is an ordered sequence of
positive integers $(c_1, \ldots, c_k)$ with $\sum c_j = n$ and
$k \geq 1$, where each part $c_j$ carries a label
$\ell_j \in \{1, \ldots, d\}$ identifying the observation channel.
\end{definition}

\begin{theorem}[Composition Inflation]
\label{thm:composition-inflation}
For $n \geq 1$ observation events across $d \geq 1$ channels, the
number of categorically distinguishable labeled trajectories is:
\begin{equation}
  T(n, d) = d\,(1 + d)^{n-1}.
  \label{eq:inflation}
\end{equation}
\end{theorem}

\begin{proof}
\textbf{Step 1.}  A composition of $n$ into $k$ positive parts
corresponds bijectively to a choice of $k-1$ dividers among $n-1$
gaps (stars-and-bars~\citep{andrews1976,graham1994}).  The number of
compositions into exactly $k$ parts is $\binom{n-1}{k-1}$.  Summing:
\begin{equation}
  \sum_{k=1}^{n} \binom{n-1}{k-1} = 2^{n-1}.
  \label{eq:total-comps}
\end{equation}

\textbf{Step 2.}  Each composition with $k$ parts receives $d^k$
independent label assignments (one label per part from $d$ channels):
\begin{equation}
  T(n,d) = \sum_{k=1}^{n} \binom{n-1}{k-1} \cdot d^k.
  \label{eq:tnd-sum}
\end{equation}

\textbf{Step 3.}  Factor out $d$ and substitute $j = k - 1$:
\begin{equation}
  T(n,d) = d \sum_{j=0}^{n-1} \binom{n-1}{j} d^j = d\,(1+d)^{n-1},
  \label{eq:tnd-binomial}
\end{equation}
by the binomial theorem.  \qed
\end{proof}

\begin{corollary}[Exponential Growth in Channels]
\label{cor:exp-growth}
For fixed $n$, $T(n,d) = \Theta(d^n)$ as $d \to \infty$.  Adding
channels inflates the trajectory space exponentially.  This inflation
occurs without any extension of the observation time window: it is
purely dimensional.
\end{corollary}

\begin{table}[h]
\centering
\caption{Values of $T(n,d) = d\,(1+d)^{n-1}$ for small $n$ and $d$.}
\label{tab:tnd}
\small
\begin{tabular}{@{}lrrrr@{}}
\toprule
$n$ \textbackslash\ $d$ & 1 & 2 & 3 & 4 \\
\midrule
1 &   1 &    2 &    3 &     4 \\
2 &   2 &    6 &   12 &    20 \\
3 &   4 &   18 &   48 &   100 \\
4 &   8 &   54 &  192 &   500 \\
5 &  16 &  162 &  768 & 2\,500 \\
6 &  32 &  486 & 3\,072 & 12\,500 \\
\bottomrule
\end{tabular}
\end{table}


% ═══════════════════════════════════════════════════════════════════════
%          PART V: THE SBS SPECIFICATION
% ═══════════════════════════════════════════════════════════════════════

\section{Systems Biology Shader Specification}
\label{sec:sbs-specification}

\subsection{Design Principles}

An SBS program is a collection of GLSL shader programs that together
constitute a biological model.  The design principles are:

\begin{enumerate}[leftmargin=*,itemsep=2pt]
\item \textbf{The shader IS the model.}  There is no separate markup
  file.  The biological model is fully specified by the shader source
  code and its uniforms.

\item \textbf{The render call IS the observation.}  Executing the
  shader pipeline on a fullscreen quad produces the observation
  texture.  No post-processing reinterpretation is required.

\item \textbf{The texture IS the result.}  The output framebuffer
  contains the measurement in S-entropy coordinates.  It is the answer,
  not a picture of the answer.

\item \textbf{No time stepping.}  The pipeline produces a single-shot
  observation.  The circuit's time invariance guarantees this shot
  captures the complete observable state.

\item \textbf{Domain specificity is localised.}  Only the partition
  observation function changes between biological domains.  All other
  pipeline stages (vertex, interference, quality, display) are
  universal.
\end{enumerate}

\subsection{Pipeline Architecture}

The SBS pipeline consists of five shader stages executed in sequence
on a fullscreen quad:

\begin{enumerate}[leftmargin=*,itemsep=2pt]
\item \textbf{Vertex Shader} (universal): Renders a fullscreen
  triangle; passes texture coordinates.
\item \textbf{Partition Observation Shader} (domain-specific):
  Computes the observation function $\Obs(\sigma)$ at each texel.
\item \textbf{Interference Shader} (universal): Computes the
  interference between two observation textures (query and reference).
\item \textbf{Quality Metric Shader} (universal): Extracts physical
  observables (partition sharpness, phase coherence, interference
  visibility, multi-resolution consistency).
\item \textbf{Display Shader} (universal): Maps the observation
  texture to a perceptually meaningful colour scheme for optional
  human inspection.
\end{enumerate}

\subsection{Universal Vertex Shader}

The vertex shader is identical for all SBS programs.  It renders a
fullscreen triangle using the vertex-ID trick, requiring no vertex
buffer.

\begin{lstlisting}[style=glsl,caption={SBS Universal Vertex Shader.}]
#version 300 es
precision highp float;

out vec2 vUV;

void main() {
    float x = -1.0 + float((gl_VertexID & 1) << 2);
    float y = -1.0 + float((gl_VertexID & 2) << 1);
    vUV = vec2(x, y) * 0.5 + 0.5;
    gl_Position = vec4(x, y, 0.0, 1.0);
}
\end{lstlisting}

\subsection{Partition Observation Shader}

The observation shader has a universal structure with a single
domain-specific hook: the \texttt{computePartitionValue} function.
This is the only component that changes between biological systems.

\begin{lstlisting}[style=glsl,caption={SBS Partition Observation Shader.}]
#version 300 es
precision highp float;

in vec2 vUV;
out vec4 fragColor;

// S-entropy coordinates (universal)
uniform float u_Sk;  // Knowledge entropy
uniform float u_St;  // Temporal entropy
uniform float u_Se;  // Evolution entropy

// Circuit state: node potentials
uniform float u_potentials[64];
// Circuit state: edge conductances
uniform float u_conductances[128];
// Network topology
uniform int u_adjacency[128];
uniform int u_numNodes;
uniform int u_numEdges;

// Domain-specific parameters
uniform float u_domainParams[16];

// ---- DOMAIN-SPECIFIC HOOK ----
vec4 computePartitionValue(
    vec2 uv,
    float Sk, float St, float Se,
    float params[16]
) {
    // Map UV to circuit node index
    int nodeIdx = int(uv.x * float(u_numNodes));
    nodeIdx = min(nodeIdx, u_numNodes - 1);

    // Read node potential (chemical potential)
    float mu = u_potentials[nodeIdx];

    // Compute partition coordinates (n,l,m,s)
    // from chemical potential via shell mapping
    float n_shell = floor(sqrt(mu * 0.5)) + 1.0;
    float l_ang = mod(mu, n_shell);
    float m_mag = mod(mu * 3.7, 2.0 * l_ang + 1.0)
                  - l_ang;
    float s_spin = sign(sin(mu * 6.28318));

    // Compute current through adjacent edges
    float J_total = 0.0;
    for (int e = 0; e < 128; e++) {
        if (e >= u_numEdges) break;
        int src = u_adjacency[e * 2];
        int dst = u_adjacency[e * 2 + 1];
        if (src == nodeIdx || dst == nodeIdx) {
            float dmu = u_potentials[src]
                      - u_potentials[dst];
            J_total += u_conductances[e] * dmu;
        }
    }

    // Triple Observation: single read,
    // three observables
    float amplitude = n_shell
                    * (1.0 - l_ang / n_shell);
    float phase = atan(m_mag, l_ang) / 6.28318;
    float depth = Se
                * length(vec2(amplitude, phase));

    // S-entropy weighting
    amplitude *= Sk;
    phase *= St;

    return vec4(
        amplitude,  // R: spectral amplitude
        phase,      // G: oscillatory phase
        depth,      // B: partition depth
        1.0         // A: confidence
    );
}
// ---- END DOMAIN-SPECIFIC HOOK ----

void main() {
    fragColor = computePartitionValue(
        vUV, u_Sk, u_St, u_Se,
        u_domainParams
    );
}
\end{lstlisting}

\subsection{Interference Shader}

The interference shader computes the superposition of two observation
textures, producing a texture whose visibility quantifies their
similarity.

\begin{lstlisting}[style=glsl,caption={SBS Universal Interference Shader.}]
#version 300 es
precision highp float;

in vec2 vUV;
out vec4 fragColor;

uniform sampler2D u_texQuery;
uniform sampler2D u_texReference;

void main() {
    vec4 q = texture(u_texQuery, vUV);
    vec4 r = texture(u_texReference, vUV);

    // Superpose as wave amplitudes
    float ampQ = q.r;
    float ampR = r.r;
    float phiQ = q.g * 6.28318;
    float phiR = r.g * 6.28318;

    // Interference: |A_q e^{i phi_q}
    //              + A_r e^{i phi_r}|^2
    float re = ampQ * cos(phiQ)
             + ampR * cos(phiR);
    float im = ampQ * sin(phiQ)
             + ampR * sin(phiR);
    float I_int = re * re + im * im;

    // Phase difference
    float dphi = phiQ - phiR;
    float cross = 2.0 * ampQ * ampR * cos(dphi);

    // Individual intensities
    float I_sum = ampQ * ampQ + ampR * ampR;

    fragColor = vec4(
        I_int,   // R: interference intensity
        cross,   // G: cross term
        I_sum,   // B: sum of intensities
        1.0      // A: weight
    );
}
\end{lstlisting}

\subsection{Quality Metric Shader}

This shader extracts four physical observables that serve as
objective, deterministic quality metrics---usable as training signals
without human labels.

\begin{lstlisting}[style=glsl,caption={SBS Universal Quality Metric Shader.}]
#version 300 es
precision highp float;

in vec2 vUV;
out vec4 fragColor;

uniform sampler2D u_texObs;
uniform sampler2D u_texInt;
uniform vec2 u_resolution;
uniform float u_sigma;

void main() {
    vec2 px = 1.0 / u_resolution;

    // --- Partition Sharpness (gradient) ---
    vec4 dx = texture(u_texObs, vUV + vec2(px.x, 0))
            - texture(u_texObs, vUV - vec2(px.x, 0));
    vec4 dy = texture(u_texObs, vUV + vec2(0, px.y))
            - texture(u_texObs, vUV - vec2(0, px.y));
    float grad = length(dx.rgb) + length(dy.rgb);

    // --- Phase Coherence ---
    vec4 ctr = texture(u_texObs, vUV);
    float phiC = ctr.g * 6.28318;
    float coherence = 0.0;
    float count = 0.0;
    for (int i = -1; i <= 1; i++) {
      for (int j = -1; j <= 1; j++) {
        if (i == 0 && j == 0) continue;
        vec2 nb = vUV + vec2(float(i),
                             float(j)) * px;
        float phiN = texture(u_texObs, nb).g
                   * 6.28318;
        coherence += cos(phiC - phiN);
        count += 1.0;
      }
    }
    coherence /= count;

    // --- Interference Visibility ---
    vec4 intTex = texture(u_texInt, vUV);
    float vis = abs(intTex.g)
              / max(intTex.b, 0.001);

    // --- Multi-Resolution Consistency ---
    vec4 blur1 = vec4(0.0);
    vec4 blur2 = vec4(0.0);
    float w1 = 0.0, w2 = 0.0;
    for (int i = -2; i <= 2; i++) {
      for (int j = -2; j <= 2; j++) {
        vec2 off = vec2(float(i), float(j));
        float d2 = dot(off, off);
        float g1 = exp(-d2
                   / (2.0 * u_sigma * u_sigma));
        float g2 = exp(-d2
                   / (8.0 * u_sigma * u_sigma));
        vec4 s = texture(u_texObs,
                         vUV + off * px);
        blur1 += s * g1; w1 += g1;
        blur2 += s * g2; w2 += g2;
      }
    }
    blur1 /= w1;
    blur2 /= w2;
    float multiRes = 1.0
        - length(blur1.rgb - blur2.rgb)
        / max(length(blur1.rgb), 0.001);

    fragColor = vec4(
        grad,       // R: partition sharpness
        coherence,  // G: phase coherence
        vis,        // B: interference visibility
        multiRes    // A: multi-res consistency
    );
}
\end{lstlisting}

\subsection{Display Shader}

For optional human inspection, the display shader maps observation
textures to perceptually meaningful colour.

\begin{lstlisting}[style=glsl,caption={SBS Universal Display Shader.}]
#version 300 es
precision highp float;

in vec2 vUV;
out vec4 fragColor;

uniform sampler2D u_texObs;
uniform float u_displayMode;
  // 0=observation, 1=interference, 2=quality

void main() {
    vec4 obs = texture(u_texObs, vUV);

    if (u_displayMode < 0.5) {
        // Observation: HSV colormap
        float h = obs.r;
        float s = 0.8;
        float v = 0.3 + 0.7 * obs.b;
        float c = v * s;
        float x = c * (1.0 - abs(
            mod(h * 6.0, 2.0) - 1.0));
        float m = v - c;
        vec3 rgb;
        float hh = h * 6.0;
        if      (hh < 1.0) rgb = vec3(c,x,0);
        else if (hh < 2.0) rgb = vec3(x,c,0);
        else if (hh < 3.0) rgb = vec3(0,c,x);
        else if (hh < 4.0) rgb = vec3(0,x,c);
        else if (hh < 5.0) rgb = vec3(x,0,c);
        else               rgb = vec3(c,0,x);
        fragColor = vec4(rgb + m, 1.0);
    } else if (u_displayMode < 1.5) {
        // Interference: blue-red diverging
        float v = obs.g * 0.5 + 0.5;
        fragColor = vec4(v, 0.2, 1.0-v, 1.0);
    } else {
        // Quality: green=good, red=bad
        float q = (obs.r + obs.g
                 + obs.b + obs.a) / 4.0;
        fragColor = vec4(1.0-q, q, 0.2, 1.0);
    }
}
\end{lstlisting}

\subsection{$\BigO(1)$-Memory Streaming Protocol}

\begin{theorem}[$\BigO(1)$ Memory Complexity]
\label{thm:o1-memory}
The SBS pipeline requires constant GPU memory independent of database
size $N$:
\begin{equation}
  \Mgpu = \Mshader + 3 \cdot W \cdot H \cdot 4
  \cdot \mathrm{sizeof(float)} = \BigO(1),
  \label{eq:gpu-memory}
\end{equation}
where $W \times H$ is the texture resolution (fixed), $\Mshader$ is
the compiled shader size (fixed), and the factor 3 accounts for the
query, reference, and interference textures.
\end{theorem}

\begin{proof}
By the Rendering--Measurement Identity (Theorem~\ref{thm:rmi}),
observation can be re-performed in constant time (one draw call).
Storing prior observations is therefore redundant: we can re-observe
any state by re-rendering.  The streaming protocol maintains three
textures (query, current reference, interference result) plus the
shader programs.  All quantities are fixed at pipeline initialisation.
When comparing a query against $N$ database entries, we stream entries
one at a time, each time re-rendering the observation and interference
shaders.  At no point does the GPU hold more than three textures plus
shaders.

For $W = H = 256$ with 4-channel 32-bit floats:
\begin{align}
  \Mgpu &= \Mshader + 3 \times 256 \times 256 \times 16 \notag \\
        &\approx 50\,\text{KB} + 3 \times 1.05\,\text{MB}
        \approx 3.2\,\text{MB}.
  \label{eq:memory-calc}
\end{align}
Including driver overhead and uniform buffers, the total is
$\sim\!13$\,MB, independent of $N$.  \qed
\end{proof}

\begin{corollary}[Integrated GPU Sufficiency]
\label{cor:integrated-gpu}
The SBS pipeline requires $\sim\!13$\,MB GPU memory and
$\sim\!1$--$2$\,TFLOPS.  Intel UHD, AMD Radeon integrated, and Apple
M-series GPUs exceed these requirements by an order of magnitude.  No
discrete GPU, no cloud, and no datacenter is needed.
\end{corollary}


% ═══════════════════════════════════════════════════════════════════════
%             PART VI: BIOLOGICAL APPLICATIONS
% ═══════════════════════════════════════════════════════════════════════

\section{Applications in Systems Biology}
\label{sec:applications}

\subsection{Cell-State Instantiation from Partial Observations}

The most fundamental SBS operation is trajectory completion
(Algorithm~\ref{alg:completion}).  Given partial omics measurements
(transcriptomics, proteomics, metabolomics), the algorithm infers the
complete molecular state.

\textbf{SBS implementation.}  The partial observations are uploaded
as shader uniforms ($\texttt{u\_potentials}[i]$ for observed nodes,
set to $-1$ for unobserved).  The observation shader implements one
iteration of KCL--KVL--backward propagation per render pass.  Multiple
render passes (ping-ponging between two framebuffers) iterate to
convergence.  The final texture contains the complete state.

The key difference from SBML-based approaches: no ODE integrator is
invoked.  The shader is not simulating the cell forward in time.  It
is propagating constraints through the circuit graph until
self-consistency is achieved.

\subsection{Disease Detection via Loop Holonomy}

\begin{definition}[Loop Holonomy]
\label{def:holonomy}
For a closed loop $\ell = (i_1, \ldots, i_k, i_1)$ in the cellular
circuit, the \emph{loop holonomy} is the parallel transport operator:
\begin{equation}
  H_\ell = \prod_{j=1}^{k} T_{i_j \to i_{j+1}},
  \label{eq:holonomy}
\end{equation}
where $T_{i \to j}$ is the transition matrix from node $i$ to node
$j$~\citep{nakahara2003,kobayashi1963}.
\end{definition}

\begin{theorem}[Health--Holonomy Identity]
\label{thm:health-holonomy}
A cellular circuit is in a healthy state if and only if the loop
holonomy is the identity for every fundamental loop:
\begin{equation}
  H_\ell = \mathrm{Id} \quad \forall\, \ell \in \pi_1(G).
  \label{eq:health}
\end{equation}
Disease corresponds to $H_\ell \neq \mathrm{Id}$ for at least one
loop.
\end{theorem}

\begin{proof}
At detailed balance (healthy equilibrium), the Wegscheider
condition~\citep{wegscheider1901} requires
$\prod_\ell k_{\mathrm{fwd}} / k_{\mathrm{rev}} = 1$ around every
cycle.  This is equivalent to $H_\ell = \mathrm{Id}$.  When disease
disrupts one or more edge conductances, the product deviates from
unity, and $H_\ell \neq \mathrm{Id}$.  The deviation
$\|H_\ell - \mathrm{Id}\|$ quantifies the severity of the
disruption.  \qed
\end{proof}

\textbf{SBS implementation.}  A dedicated SBS shader computes the
holonomy for each fundamental loop of the circuit by accumulating
transition matrices along the loop path.  The output texture encodes
$\|H_\ell - \mathrm{Id}\|$ at each texel, with healthy regions
appearing as zero (black) and diseased regions as nonzero (coloured).

\subsection{Coherence as Health Index}

\begin{definition}[Cellular Coherence Index]
\label{def:coherence}
The \emph{cellular coherence index} is the Kuramoto order
parameter~\citep{kuramoto1975,strogatz2000,acebron2005} of the
cellular oscillator ensemble:
\begin{equation}
  R = \left| \frac{1}{N} \sum_{j=1}^{N} e^{i\varphi_j} \right|,
  \label{eq:kuramoto}
\end{equation}
where $\varphi_j$ is the oscillatory phase of the $j$-th molecular
species.
\end{definition}

\begin{theorem}[Coherence--Health Identity]
\label{thm:coherence-health}
The multi-channel interference visibility $V_{\mathrm{cell}}$ of the
SBS observation texture equals the Kuramoto order parameter $R$:
\begin{equation}
  V_{\mathrm{cell}} = R.
  \label{eq:coherence-health}
\end{equation}
Healthy cells have $R > 0.7$ (synchronised oscillators); diseased
cells have $R < 0.3$ (desynchronised).
\end{theorem}

\begin{proof}
The phase $\varphi_j$ accumulated along observation channel $j$ is
the G-channel value of the observation texture at the corresponding
texel (equation~\ref{eq:rgba-observables}).  The interference
visibility (computed by the interference shader) is
$V = |I_{\max} - I_{\min}| / (I_{\max} + I_{\min})$, which equals
$|\langle e^{i\varphi_j} \rangle|$ when the observation channels
sample the molecular ensemble uniformly.  By
definition~\eqref{eq:kuramoto}, this is $R$.  \qed
\end{proof}

\subsection{Therapeutic Effect Prediction}

\begin{definition}[Drug Perturbation]
\label{def:drug}
A drug is a sparse perturbation $\boldsymbol{\eta} \in [0,1]^R$ of
the edge conductances: $\Gcond_{ij} \to \Gcond_{ij}(1 - \eta_{ij})$,
where $\eta_{ij} = 0$ means no effect on edge $(i,j)$ and
$\eta_{ij} = 1$ means complete inhibition.
\end{definition}

\begin{theorem}[Optimal Drug Design as $\ell_1$ Minimisation]
\label{thm:drug-design}
Given a diseased cellular circuit with loop holonomies
$\{H_\ell \neq \mathrm{Id}\}$, the optimal drug perturbation
$\boldsymbol{\eta}^*$ minimises the therapeutic score:
\begin{equation}
  \boldsymbol{\eta}^* = \arg\min_{\boldsymbol{\eta}}
  \left\{
    \sum_\ell \|H_\ell(\boldsymbol{\eta}) - \mathrm{Id}\|^2
    + \lambda \|\boldsymbol{\eta}\|_1
  \right\},
  \label{eq:drug-opt}
\end{equation}
where $\lambda$ is the side-effect penalty.  This is a convex
program~\citep{boyd2004} solvable in polynomial time.
\end{theorem}

\begin{proof}
The holonomy $H_\ell(\boldsymbol{\eta})$ depends continuously on the
conductances, which are affine in $\boldsymbol{\eta}$.  The quadratic
loop-holonomy term is convex in the conductances (sum of squared
norms).  The $\ell_1$ penalty is convex by definition.  The sum of
convex functions is convex.  Standard convex programming
methods~\citep{boyd2004} solve this in polynomial time, yielding a
sparse $\boldsymbol{\eta}^*$ (most edges unperturbed) that restores
$H_\ell \approx \mathrm{Id}$ for all diseased loops.  \qed
\end{proof}

\textbf{SBS implementation.}  The drug perturbation
$\boldsymbol{\eta}$ is uploaded as the
$\texttt{u\_conductances}$ uniform array (with perturbed values).
Two render passes are executed: pre-treatment and post-treatment.  The
therapeutic score is computed by the quality metric shader as the
difference in interference visibility between the two passes.  The
full prediction---patient image to treatment decision---completes in
under 100\,ms on a laptop GPU.

\subsection{Variance--Free Energy Identity}

\begin{theorem}[Variance--Free Energy Correspondence]
\label{thm:variance-free-energy}
The phase variance $\sigma^2(\varphi)$ of the cellular oscillator
ensemble is proportional to the partition free energy:
\begin{equation}
  F = \kb T \cdot \sigma^2(\varphi).
  \label{eq:variance-free-energy}
\end{equation}
\end{theorem}

\begin{proof}
The free energy is $F = -\kb T \ln Z$, where $Z$ is the partition
function.  For oscillators with phases $\varphi_j$ distributed around
mean $\bar{\varphi}$, the partition function is
$Z = \int_0^{2\pi} e^{-\beta(\varphi - \bar{\varphi})^2/2}\,d\varphi
\propto 1/\sqrt{\sigma^2}$ in the Gaussian approximation.  Then
$F = -\kb T \ln(1/\sqrt{\sigma^2}) = \tfrac{1}{2}\kb T \ln \sigma^2$,
which is monotonically increasing in $\sigma^2$.  To leading order,
$F \propto \kb T \cdot \sigma^2(\varphi)$ for small deviations.
A decrease in phase variance corresponds to a decrease in free
energy---exactly what a successful therapy achieves: re-synchronising
the cellular oscillators reduces both variance and free energy.  \qed
\end{proof}


\subsection{GPU-Supervised Training}

\begin{theorem}[Label-Free Training Signal]
\label{thm:gpu-supervised}
The four quality metrics extracted by the SBS quality metric shader
(partition sharpness $P_{\mathrm{sharp}}$, phase coherence
$C_{\mathrm{phase}}$, interference visibility $V$, multi-resolution
consistency $M_{\mathrm{res}}$) are objective physical observables
that provide a training loss without human labels:
\begin{equation}
  \mathcal{L}_{\mathrm{GPU}} = w_1(1 - P_{\mathrm{sharp}})
  + w_2(1 - C_{\mathrm{phase}}) + w_3(1 - V)
  + w_4(1 - M_{\mathrm{res}}).
  \label{eq:gpu-loss}
\end{equation}
Training proceeds even when the number of human-labelled examples is
zero.
\end{theorem}

\begin{proof}
Each metric is a deterministic function of the observation texture,
computed by the quality metric shader.  By the Rendering--Measurement
Identity (Theorem~\ref{thm:rmi}), these are physical measurements,
not statistical estimates.  They are: (i)~deterministic (same input
$\to$ same output); (ii)~differentiable (smooth functions of texture
values); (iii)~informative (each captures an independent aspect of
observation quality).  The gradient
$\nabla_\theta \mathcal{L}_{\mathrm{GPU}}$ with respect to model
parameters $\theta$ can be computed by backpropagation through the
shader pipeline.  No human label $y$ appears in the
loss.  \qed
\end{proof}


% ═══════════════════════════════════════════════════════════════════════
%           PART VII: COMPARISON, VALIDATION, AND CONCLUSIONS
% ═══════════════════════════════════════════════════════════════════════

\section{Comparison with Existing Formalisms}
\label{sec:comparison}

\begin{table*}[t]
\centering
\caption{Comparison of systems biology formalisms.}
\label{tab:comparison}
\small
\begin{tabular}{@{}lccccc@{}}
\toprule
\textbf{Property} &
  \rotatebox{70}{\textbf{SBML}} &
  \rotatebox{70}{\textbf{CellML}} &
  \rotatebox{70}{\textbf{BioPAX}} &
  \rotatebox{70}{\textbf{Petri nets}} &
  \rotatebox{70}{\textbf{SBS}} \\
\midrule
Model representation     & Markup  & Markup  & Markup  & Graph   & Shader \\
Computation method       & ODE solver & ODE solver & N/A & Token sim. & Render \\
Observation = computation & No & No & No & No & \textbf{Yes} \\
Time stepping required   & Yes & Yes & N/A & Yes & \textbf{No} \\
Memory scaling           & $\BigO(N)$ & $\BigO(N)$ & $\BigO(N)$ & $\BigO(N)$ & $\BigO(1)$ \\
GPU native               & No & No & No & No & \textbf{Yes} \\
Label-free training      & No & No & No & No & \textbf{Yes} \\
Disease detection        & Via simulation & Via simulation & N/A & Via reachability & \textbf{Via holonomy} \\
Integrated GPU sufficient & N/A & N/A & N/A & N/A & \textbf{Yes} \\
\bottomrule
\end{tabular}
\end{table*}

The essential difference between SBS and all prior formalisms is
ontological.  SBML, CellML, and BioPAX describe biological models as
data structures to be processed by external algorithms.  The model and
the computation are separate entities.  In SBS, the model \emph{is}
the computation: the shader program is simultaneously the model
specification and the observation apparatus.  This is not a
convenience---it is a consequence of the Rendering--Measurement
Identity (Theorem~\ref{thm:rmi}).

A second fundamental difference is temporal.  All markup-based
formalisms require a time-stepping algorithm (ODE integrator,
Gillespie sampler, token simulator) that advances the system state
through a sequence of time points.  SBS requires no time stepping.
The time invariance of backward trajectories
(Theorem~\ref{thm:time-invariance}) guarantees that a single render
pass captures the complete observable state.

A third difference is memory scaling.  Markup-based approaches require
$\BigO(N)$ memory to store the model state, where $N$ is the number of
species.  SBS requires $\BigO(1)$ GPU memory
(Theorem~\ref{thm:o1-memory}), because the observation texture has
fixed resolution independent of model size.


\section{Validation Design}
\label{sec:validation}

We propose three classes of validation experiments.

\subsection{Experiment 1: Shell Capacity Validation}

\textbf{Prediction.}  The partition shell capacity $C(n) = 2n^2$
(Theorem~\ref{thm:shell-capacity}) predicts the electron
configuration of atoms without any quantum-mechanical input.

\textbf{Protocol.}  For elements with atomic number $Z \leq 64$,
compute $C(n)$ for $n = 1, 2, \ldots$ and compare the cumulative
capacity $\sum_{n'=1}^{n} C(n') = \frac{2}{3}n(n+1)(2n+1)/6$ against
the experimentally observed electron shell filling from the NIST
Atomic Spectra Database~\citep{nist_asd}.

\textbf{Success criterion.}  Shell capacities match exactly
($C(1) = 2$, $C(2) = 8$, $C(3) = 18$, $C(4) = 32$) for the first
four shells.

\subsection{Experiment 2: Spectral Reconstruction}

\textbf{Prediction.}  The SBS partition observation shader, when
parameterised with the S-entropy coordinates of a known atomic species,
produces a spectral image whose peak positions match NIST reference
wavelengths~\citep{nist_asd,herzberg1944}.

\textbf{Protocol.}  Encode the S-entropy coordinates of hydrogen (H),
molecular hydrogen (H$_2$), and water (H$_2$O) as shader uniforms.
Execute the SBS observation pipeline.  Extract peak positions from the
R-channel of the output texture.  Compare against NIST reference
values.

\textbf{Success criterion.}  Mean relative error $< 0.1\%$ across at
least 20 spectral lines per species.

\subsection{Experiment 3: Disease Detection via Coherence}

\textbf{Prediction.}  The cellular coherence index $R$
(Theorem~\ref{thm:coherence-health}) separates healthy ($R > 0.7$)
from diseased ($R < 0.3$) cells from a single observation texture.

\textbf{Protocol.}  Construct SBS circuit models for: (i)~a healthy
glycolytic network (10 nodes, 15 edges, NIST/KEGG rate constants);
(ii)~the same network with one edge conductance reduced by 90\%
(simulating enzyme deficiency).  Execute the SBS pipeline for both.
Compute $R$ from the interference visibility of the quality metric
shader.

\textbf{Success criterion.}  $R_{\mathrm{healthy}} > 0.7$ and
$R_{\mathrm{diseased}} < 0.3$ with separation $> 0.4$.

\subsection{Experiment 4: Therapeutic Score Validation}

\textbf{Prediction.}  The $\ell_1$-optimal drug perturbation
$\boldsymbol{\eta}^*$ (Theorem~\ref{thm:drug-design}) restores loop
holonomy to within $\|H_\ell - \mathrm{Id}\| < 0.01$ for the diseased
network from Experiment~3.

\textbf{Protocol.}  Solve the convex program~\eqref{eq:drug-opt} for
the diseased glycolytic network.  Upload $\boldsymbol{\eta}^*$ as
modified conductances.  Re-render the SBS pipeline.  Measure
post-treatment coherence $R_{\mathrm{post}}$ and holonomy residual.

\textbf{Success criterion.}  $R_{\mathrm{post}} > 0.7$;
$\|\boldsymbol{\eta}^*\|_0 \leq 3$ (at most 3 edges perturbed);
total pipeline time $< 100$\,ms on an integrated GPU.


\section{Discussion}
\label{sec:discussion}

\subsection{What SBS Is Not}

SBS is not a faster way to solve ODEs on GPUs.  That approach---porting
CVODE to CUDA---has been explored extensively~\citep{owens2008,
nickolls2008} and remains fundamentally limited by the simulation
problem (Proposition~\ref{prop:sim-equiv}).  SBS eliminates the ODE
step entirely.  The shader does not integrate equations; it propagates
constraints through a circuit graph.

SBS is not a visualisation framework.  The Rendering--Measurement
Identity (Theorem~\ref{thm:rmi}) proves that the rendering \emph{is}
the measurement.  The texture is the result, not a picture of the
result.  The display shader (Listing~5) is optional and exists only
for human convenience; the computation is complete after the
observation shader (Listing~2).

\subsection{Relationship to the von Neumann Bottleneck}

The von Neumann bottleneck~\citep{backus1978}---the bandwidth
limitation between processor and memory---is the root cause of the
$\BigO(N)$ memory scaling of conventional approaches.  SBS eliminates
this bottleneck by proving that storage is redundant
(Theorem~\ref{thm:o1-memory}): if observation can be re-performed in
constant time, there is no need to store results.  The GPU's
massively parallel evaluation of the fragment shader across all texels
simultaneously replaces the sequential fetch-decode-execute cycle of
the CPU.

\subsection{Universality Beyond Biology}

The SBS architecture is domain-agnostic.  The only domain-specific
component is the \texttt{computePartitionValue} function in the
observation shader.  All other components---vertex shader, interference
shader, quality metric shader, display shader, streaming protocol,
memory analysis---are universal.  Replacing the biological circuit
model with a financial time-series model, a climate model, or a
materials-science model requires changing only one function.  The
Rendering--Measurement Identity holds for any domain, because it is a
theorem about the relationship between functions and their
evaluations, not about biology.

\subsection{Limitations}

\textbf{Circuit topology.}  The current formalism requires the
reaction network topology to be known.  For poorly characterised
organisms, this limits applicability.  However, the circuit structure
can be inferred from partial observations via the trajectory completion
algorithm (Algorithm~\ref{alg:completion}), and databases such as
KEGG and Reactome provide curated topologies for model organisms.

\textbf{Texture resolution.}  The observation texture has fixed
resolution $W \times H$.  For circuits with more than $W \times H$
nodes, tiling or hierarchical observation strategies are required.
At $256 \times 256 = 65{,}536$ texels, this accommodates the largest
curated metabolic networks (human Recon3D: $\sim$5{,}000 metabolites,
$\sim$13{,}000 reactions).

\textbf{Fuzzy initialisation.}  The fuzzy membership functions
(Definition~\ref{def:fuzzy}) require bounds on measurement
uncertainty.  In practice, these are available from instrument
specifications and replicate measurements.

\textbf{Oscillator jitter.}  Timing cells
(Definition~\ref{def:timing-cell}) must have width exceeding the
oscillator jitter bound to avoid false cell assignments.  For commodity
crystal oscillators ($\epsilon \approx 10$\,ppm), minimum cell widths
are $\sim\!10$\,ns, which is adequate for all biological
timescales.

\subsection{Relationship to Temporal Programming}

The timing-cell formalism (\S\ref{sec:temporal-observation}) embeds a
complete programming model within SBS.  A \emph{timing cell} is the
biological analog of a Tempus cell declaration: a measurable interval
in $\DP$-space with a pre-compiled action.  The cell partition is
the biological analog of a cell registry stored in ROM.  The
Composition Inflation Theorem
(Theorem~\ref{thm:composition-inflation}) provides the state-space
capacity.

This connection runs deeper than analogy.  The COMPILE/EXECUTE phase
separation of temporal programming---where trajectory accumulation and
action dispatch are mutually exclusive---maps onto the SBS pipeline:
the observation shader accumulates (reads state), and the interference
shader dispatches (computes the comparison).  These are separate render
passes that never execute simultaneously.

\subsection{Structural Security}

Because the SBS pipeline processes only timing deviations and shader
uniforms---not parsed payloads---the classical attack surfaces of
computational biology pipelines (buffer overflow in file parsers,
injection in query engines, man-in-the-middle on data transfers) are
architecturally absent.  The shader programs are compiled at deploy
time and stored as immutable GPU binaries.  No runtime code generation,
JIT compilation, or dynamic loading occurs.  The attack surface is the
finite set of uniform values, which are bounded by construction
($[0,1]^3$ for S-entropy coordinates).


\section{Conclusion}
\label{sec:conclusion}

We have introduced Systems Biology Shaders (SBS), a formalism that
replaces the markup-then-simulate paradigm of SBML and CellML with a
shader-is-the-model paradigm in which biological systems are specified
as GPU fragment shaders.

The theoretical foundation rests on three independently derived
results.  The \emph{Oscillatory Necessity Theorem} proves that every
persistent bounded system oscillates---this is not a modelling
choice but a mathematical consequence of boundedness.  The
\emph{Triple Equivalence Theorem} proves that oscillation, categorical
distinction, and partition operation are the same mathematical
object---connected by explicit bijections with identical state counts
and entropies.  The \emph{Rendering--Measurement Identity} proves that
a fragment shader evaluating a partition observation function
\emph{is} performing a physical measurement---the texture is the
result, not a picture of it.

From these foundations we derived a cellular circuit model with exact
Kirchhoff analogs, time-invariant backward trajectories that serve as
categorical addresses, S-entropy coordinates that map directly to GPU
textures, and a Composition Inflation Theorem showing that
multi-channel observation provides $T(n,d) = d\,(1+d)^{n-1}$
distinguishable trajectories through dimensional progression without
temporal evolution.

The practical consequence is a five-stage GLSL shader pipeline
(vertex, observation, interference, quality, display) that runs on any
WebGL\,2 browser with $\sim\!13$\,MB GPU memory, requires no time
stepping, and produces cell-state observations, disease diagnoses, and
therapeutic predictions in under 100\,ms on integrated laptop GPUs.

SBS does not accelerate existing computations.  It eliminates them.
The biological model is the shader.  The observation is the render
call.  The result is the texture.  Nothing else is needed.


% ═══════════════════════════════════════════════════════════════════════
%                         REFERENCES
% ═══════════════════════════════════════════════════════════════════════

\bibliographystyle{plainnat}
\bibliography{references}


% ═══════════════════════════════════════════════════════════════════════
%                         APPENDICES
% ═══════════════════════════════════════════════════════════════════════

\appendix

\section{SBS File Format Specification}
\label{app:format}

An SBS file is a directory containing:

\begin{itemize}[leftmargin=*,nosep]
\item \texttt{manifest.json}: Metadata (organism, network ID, version,
  S-entropy ranges, number of nodes and edges).
\item \texttt{vertex.glsl}: Vertex shader (typically the universal
  shader from Listing~1).
\item \texttt{observation.glsl}: Partition observation shader with the
  domain-specific \texttt{computePartitionValue} function.
\item \texttt{interference.glsl}: Interference shader (typically
  universal, Listing~3).
\item \texttt{quality.glsl}: Quality metric shader (typically
  universal, Listing~4).
\item \texttt{display.glsl}: Display shader (typically universal,
  Listing~5).
\item \texttt{circuit.json}: Circuit topology (adjacency list, initial
  potentials, conductances).
\item \texttt{entropy.json}: S-entropy coordinate ranges and
  normalisation parameters.
\end{itemize}

The manifest specifies which shader files are universal and which are
domain-specific, enabling tools to validate and optimise the pipeline
without parsing biological content.

\section{Proof of Convergence Rate}
\label{app:convergence-rate}

We provide the explicit contraction constant for
Theorem~\ref{thm:convergence}.

Let $\lambda_2$ be the second-smallest eigenvalue of the normalised
graph Laplacian $\mathcal{L} = I - D^{-1/2}AD^{-1/2}$, where $A$ is
the adjacency matrix and $D$ is the degree matrix.  The KCL averaging
step contracts by a factor $(1 - \lambda_2)$.  The KVL loop closure
step eliminates the component of the error lying in the cycle space,
which for a connected graph with $\beta_1 = R - N + 1$ independent
cycles projects onto a subspace of codimension $\beta_1$.  The Viterbi
backward pass is a projection onto a finite set of paths, which is
a non-expansive map.

The composite contraction constant is:
\begin{equation}
  \rho \leq (1 - \lambda_2) \cdot
  \sqrt{1 - \frac{\beta_1}{R}} \cdot 1
  = (1 - \lambda_2)\sqrt{1 - \frac{\beta_1}{R}}.
  \label{eq:contraction-constant}
\end{equation}

For a well-connected metabolic network with $\lambda_2 \approx 0.1$
and $\beta_1/R \approx 0.3$, the contraction constant is
$\rho \approx 0.9 \times 0.84 \approx 0.75$, yielding convergence to
$\epsilon$-accuracy in $\lceil \log(\epsilon) / \log(0.75) \rceil
\approx 25$ iterations for $\epsilon = 10^{-8}$.


\section{Computational Complexity Summary}
\label{app:complexity}

\begin{table}[H]
\centering
\caption{Per-observation computational complexity.}
\label{tab:complexity}
\small
\begin{tabular}{@{}lcc@{}}
\toprule
\textbf{Operation} & \textbf{SBS} & \textbf{SBML+solver} \\
\midrule
Model parsing          & $\BigO(1)$ (compiled) & $\BigO(N)$ \\
State observation      & $\BigO(1)$ (render)   & $\BigO(N \cdot T)$ \\
Database comparison    & $\BigO(1)$ (per item) & $\BigO(N \cdot d)$ \\
Memory (GPU)           & $\BigO(1)$            & $\BigO(N)$ \\
Disease detection      & $\BigO(|\pi_1(G)|)$   & $\BigO(N \cdot T)$ \\
Drug optimisation      & $\BigO(R \cdot |\pi_1(G)|)$ & $\BigO(N \cdot T \cdot R)$ \\
\bottomrule
\multicolumn{3}{l}{\footnotesize $N$ = species, $R$ = reactions,
$T$ = time steps, $d$ = embedding dim.}
\end{tabular}
\end{table}


\end{document}
